

tardará en comenzar a recibir efectivo propio de la inversión. Para ello tomaremos tantos factores como sean necesarios al proyecto, además restamos el porcentaje de las ventas de contado por ser el dinero que ya estamos recibiendo.

De acuerdo a la estructura del plan de inversión para la producción del servicio mostrado en la Figura 10.19, necesitamos un monto total de S/21,154.29, de los cuales el 69.07\%, es decir aproximadamente $\mathrm{S} / 25,125.67$ corresponde a fondos propios de la sociedad y el $30.93 \%$ restante debe ser cubierto por el sistema financiero mediante una Operación de Leasing.

El Leasing debe ser entendido como una operación de carácter financiero que consiste en la adquisición por parte la entidad financiera del equipamiento y la infraestructura y la cesión simultánea de uso durante un período de tiempo determinado que en nuestro caso es de seis años, a un precío que se distribuye en cuotas periódicas.
\subsubsection{ Balance de Apertura}

El Balance mostrado en la Figura 10.20 se elabora en base al Plan de Inversión de la Figura 10.19; caja y bancos es igual al capital de trabajo necesitado para iniciar operaciones. El pasivo es igual al saldo del financiamiento: se toma como pasivo a corto plazo el importe del primer pago a capital según la forma de pago escogida; cuota nivelada o cuota de saldos insolutos. En este proyecto se decidió por la cuota nivelada Figura 10.13. El pasivo a largo plazo sería la diferencia entre lo que se tiene que pagar en este año y el importe total del financiamiento. El patrimonio será igual al total de activo fijo, es decir, lo que hemos invertido en el proyecto.
